\chapter{Expert Scenario Evaluation Sheet}
\label{sec:appendix_expert_eval}
 
This appendix contains the complete expert evaluation instrument administered during the scenario
validation phase of the study. Section~\ref{sec:app_scoring} presents the expert's criterion-level
ratings and qualitative comments for each of the ten evaluated scenarios. 
Section~\ref{sec:app_scenarios} presents the full scenario descriptions submitted to the expert,
including all eleven documentation dimensions used for evaluation.
 
%-----------------------------------------------------------------------------
% SECTION 1: EXPERT SCORING SHEET
%-----------------------------------------------------------------------------
 
\section{Expert Criterion Ratings and Comments}
\label{sec:app_scoring}
 
The table below presents the expert's ratings on the 3-point Likert scale
(0 = Not Suitable, 1 = Needs Revision, 2 = Suitable) across all six evaluation
criteria, alongside any qualitative comments provided per scenario.
Abbreviations: \textbf{GM} = Grice's Maxims Alignment; \textbf{YA} = Relevance
to Young Adult Contexts; \textbf{CA} = Cultural Appropriateness; \textbf{LV} =
Learning Value; \textbf{CP} = Complexity \& Progression; \textbf{CS} = Clarity
\& Structure.
 
\begin{landscape}
 
\begin{longtable}{>{\bfseries}p{0.4cm} p{3.4cm} >{\centering\arraybackslash}p{0.6cm} >{\centering\arraybackslash}p{0.6cm} >{\centering\arraybackslash}p{0.6cm} >{\centering\arraybackslash}p{0.6cm} >{\centering\arraybackslash}p{0.6cm} >{\centering\arraybackslash}p{0.6cm} p{6.8cm}}
\caption{Expert Evaluation Scoring Sheet with Qualitative Comments}
\label{tab:app_expert_scores} \\
 
\rowcolor{headerblue}
\textcolor{white}{\textbf{\#}} &
\textcolor{white}{\textbf{Scenario Title}} &
\textcolor{white}{\textbf{GM}} &
\textcolor{white}{\textbf{YA}} &
\textcolor{white}{\textbf{CA}} &
\textcolor{white}{\textbf{LV}} &
\textcolor{white}{\textbf{CP}} &
\textcolor{white}{\textbf{CS}} &
\textcolor{white}{\textbf{Expert Comments}} \\
\endfirsthead
 
\multicolumn{9}{c}{\tablename\ \thetable{} -- continued} \\
\rowcolor{headerblue}
\textcolor{white}{\textbf{\#}} &
\textcolor{white}{\textbf{Scenario Title}} &
\textcolor{white}{\textbf{GM}} &
\textcolor{white}{\textbf{YA}} &
\textcolor{white}{\textbf{CA}} &
\textcolor{white}{\textbf{LV}} &
\textcolor{white}{\textbf{CP}} &
\textcolor{white}{\textbf{CS}} &
\textcolor{white}{\textbf{Expert Comments}} \\
\endhead
 
\midrule
1 & The Course Change Dilemma            & 2 & 2 & 2 & 2 & 2 & 2 & --- \\
\midrule
2 & The ``Pauwi Na'' Curfew              & 2 & 2 & 2 & 2 & 2 & 2 & --- \\
\midrule
3 & The Ambiguous ``Kayo Na Ba?''        & 2 & 2 & 2 & 1 & 0 & 0 &
There are two things happening in this scenario: \textit{hiya} between the ``couple''; and
\textit{hiya} between the couple and other people. The need for clarification regarding the
relationship status accordingly has two levels. It is recommended to limit the conflict to
between the couple. A rich conversation can already ensue in that single dynamic. \\
\midrule
4 & The ``Sama Ka'' Project Group        & 2 & 2 & 2 & 2 & 2 & 2 &
Consider a variant involving balancing academic and org/extra-curricular work. \\
\midrule
5 & The Family Business ``Opportunity''  & 2 & 2 & 2 & 2 & 2 & 2 & --- \\
\midrule
6 & The ``Wala Naman Akong Pera'' Friend & 2 & 2 & 2 & 2 & 2 & 2 &
An additional scenario to consider: a friend offering an insurance package which the user
cannot afford. The tension between being a good friend and maintaining financial autonomy
could be a strong pragmatic conflict. (\textit{Implemented as Scenario 11.}) \\
\midrule
7 & The ``Okay Lang Ako'' Partner        & 2 & 2 & 2 & 2 & 2 & 2 & --- \\
\midrule
8 & The Political Family Dinner          & 2 & 2 & 2 & 2 & 2 & 2 & --- \\
\midrule
9 & The Micromanaging Boss               & 2 & 2 & 2 & 2 & 2 & 2 & --- \\
\midrule
10 & The ``Chismis'' Lunch Break         & 2 & 2 & 2 & 2 & 2 & 2 & --- \\
\bottomrule
\end{longtable}
 
\end{landscape}
 
\noindent \textbf{Revisions Implemented Following Expert Feedback.} Based on the expert's
comments, the following changes were made prior to system deployment:
 
\begin{itemize}
  \item \textbf{Scenario 3} (The Ambiguous ``Kayo Na Ba?''): The conflict was narrowed to a
  strictly one-on-one dynamic between the couple only, removing the secondary layer involving
  external parties. This resolved the Complexity \& Progression (0) and Clarity \& Structure (0)
  flags by bounding the scenario's pragmatic scope.
 
  \item \textbf{Scenario 4} (The ``Sama Ka'' Project Group): The expert's suggestion regarding
  balancing academic and organizational commitments was noted. The scenario's existing conflict
  structure was retained as it already addressed in-group academic pressure; the suggestion was
  logged for future scenario expansion.
 
  \item \textbf{Scenario 6 / Scenario 11}: The expert's suggestion of a friend-as-insurance-agent
  scenario was adopted as a new eleventh scenario (\textit{The ``Insurance Agent'' Friend}),
  which focuses on navigating \textit{pakikisama} and financial autonomy when a close friend
  is the sales agent.
\end{itemize}
 
%-----------------------------------------------------------------------------
% SECTION 2: FULL SCENARIO DESCRIPTIONS
%-----------------------------------------------------------------------------
 
\section{Full Scenario Descriptions Submitted for Expert Evaluation}
\label{sec:app_scenarios}
 
The following tables present the complete documentation submitted to the expert for each of the
ten evaluated scenarios, structured across the eleven evaluation dimensions. Scenario 11 was added
after the expert validation phase based on the expert's suggestion (see above) and therefore
was not part of the formal evaluation instrument.
 
%--- Scenario 1 ---%
\subsection*{Scenario 1 --- The Course Change Dilemma}
 
\begin{longtable}{p{4.5cm} p{11cm}}
\caption*{Scenario 1: The Course Change Dilemma} \\
\toprule
\textbf{Dimension} & \textbf{Description} \\
\midrule
\endhead
Recommended Age Range          & 18--19 \\
\midrule
Thematic Classification        & Identity vs.\ Expectation \\
\midrule
Developmental Task (Havighurst)& Preparing for an economic career \\
\midrule
User Role                      & First year university student \\
\midrule
Chatbot Role                   & Parent who is financing the user's education \\
\midrule
Pragmatic Conflict             & You are on a video call with your parent, feeling anxious but determined.
                                 After struggling for a semester in a pre-med program, you have decided to
                                 shift to the fine arts course you have always dreamed of. Your parent, a
                                 well-respected doctor, listens without interruption. When you finish, their
                                 voice is quiet but carries the weight of disappointment: ``\textit{Anak, we
                                 have sacrificed so much for this opportunity. All your titos and titas, our
                                 whole family, is so proud of our future doctor. Are you really going to
                                 throw that all away\ldots to paint?}'' \\
\midrule
Key Gricean Maxims             & Quality, Relation \\
\midrule
Key Filipino Values            & \textit{Utang na Loob}, \textit{Paggalang} \\
\midrule
Expected Pragmatic Skill       & Assertive empathy. The user must articulate their personal choice and
                                 passion clearly while acknowledging the family's sacrifices and expressing
                                 gratitude, balancing the assertion of autonomy with the cultural value of
                                 \textit{Utang na Loob}. \\
\bottomrule
\end{longtable}
 
%--- Scenario 2 ---%
\subsection*{Scenario 2 --- The ``Pauwi Na'' Curfew}
 
\begin{longtable}{p{4.5cm} p{11cm}}
\caption*{Scenario 2: The ``Pauwi Na'' Curfew} \\
\toprule
\textbf{Dimension} & \textbf{Description} \\
\midrule
\endhead
Recommended Age Range          & 22--24 \\
\midrule
Thematic Classification        & Autonomy vs.\ Dependence \\
\midrule
Developmental Task (Havighurst)& Achieving emotional independence of parents and other adults \\
\midrule
User Role                      & A young professional who still lives with their parents \\
\midrule
Chatbot Role                   & A loving but overprotective mother \\
\midrule
Pragmatic Conflict             & It's 10:30 PM, and you're at a celebratory dinner with your team after
                                 successfully launching a major project. Your phone lights up. It's a message
                                 from your Mama. She has sent you a photo of the family dog sitting by the
                                 front door with the caption: ``\textit{Hachi is still waiting for you.
                                 It's getting late, baka mapano ka.}'' \\
\midrule
Key Gricean Maxims             & Manner, Relation \\
\midrule
Key Filipino Values            & \textit{Paggalang} \\
\midrule
Expected Pragmatic Skill       & Reframing the conversation from obedience to responsibility. The user
                                 must reassure their mother of their safety and maturity without being
                                 dismissive, shifting the parent-child dynamic toward one of adult-to-adult
                                 respect. \\
\bottomrule
\end{longtable}
 
%--- Scenario 3 ---%
\subsection*{Scenario 3 --- The Ambiguous ``Kayo Na Ba?''}
 
\begin{longtable}{p{4.5cm} p{11cm}}
\caption*{Scenario 3: The Ambiguous ``Kayo Na Ba?'' \textnormal{(\textit{Revised after expert feedback})}} \\
\toprule
\textbf{Dimension} & \textbf{Description} \\
\midrule
\endhead
Recommended Age Range          & 19--21 \\
\midrule
Thematic Classification        & Ambiguity vs.\ Definition \\
\midrule
Developmental Task (Havighurst)& Achieving new and more mature relations with age-mates of both sexes \\
\midrule
User Role                      & A college student in a ``talking stage'' \\
\midrule
Chatbot Role                   & The person the user has been exclusively dating for months without a
                                 label \\
\midrule
Pragmatic Conflict             & You and this person do everything a couple does: go on dates, meet each
                                 other's friends, and are physically affectionate. But the relationship
                                 remains undefined. Tonight, while walking through campus, they abruptly
                                 let go of your hand. Later, they send a text that feels too casual:
                                 ``\textit{Had fun tonight! :) See you tomorrow?}'' \\
\midrule
Key Gricean Maxims             & Quantity, Quality \\
\midrule
Key Filipino Values            & \textit{Pakikisama}, \textit{Hiya}, \textit{Kapwa} \\
\midrule
Expected Pragmatic Skill       & Clarification through tactful questioning. The user must navigate the
                                 cultural fear of rejection (\textit{Hiya}) to gently express their need
                                 for clarity about the relationship's status, turning an ambiguous situation
                                 into a defined one. \\
\midrule
Post-Expert Revision Note      & Conflict limited to the couple's one-on-one dynamic only; the secondary
                                 layer involving \textit{hiya} with external parties was removed. \\
\bottomrule
\end{longtable}
 
%--- Scenario 4 ---%
\subsection*{Scenario 4 --- The ``Sama Ka'' Project Group}
 
\begin{longtable}{p{4.5cm} p{11cm}}
\caption*{Scenario 4: The ``Sama Ka'' Project Group} \\
\toprule
\textbf{Dimension} & \textbf{Description} \\
\midrule
\endhead
Recommended Age Range          & 16--18 \\
\midrule
Thematic Classification        & Group vs.\ Self \\
\midrule
Developmental Task (Havighurst)& Acquiring a set of values and an ethical system as a guide to behavior \\
\midrule
User Role                      & A senior high school student in a group project \\
\midrule
Chatbot Role                   & The popular and influential leader of the user's project group \\
\midrule
Pragmatic Conflict             & Your final paper is due tomorrow, and your group is far from finished.
                                 The group leader, one of the most popular students in your class, sends
                                 a message to the group chat: ``\textit{Sorry guys, wala pa akong nagawa,
                                 busy kasi sa org eh. But relax, may AI naman tayo. The prof will never
                                 notice. Just don't tell anyone. We're all in this together, right?
                                 Sama kayo, di ba?}'' \\
\midrule
Key Gricean Maxims             & Quality \\
\midrule
Key Filipino Values            & \textit{Pakikisama}, \textit{Hiya} \\
\midrule
Expected Pragmatic Skill       & Intervening respectfully and de-escalating conflict. The user must voice
                                 their ethical concerns and refuse to plagiarize, risking social friction
                                 (\textit{Hiya}) while persuading the group toward a better solution and
                                 maintaining harmony (\textit{Pakikisama}). \\
\bottomrule
\end{longtable}
 
%--- Scenario 5 ---%
\subsection*{Scenario 5 --- The Family Business ``Opportunity''}
 
\begin{longtable}{p{4.5cm} p{11cm}}
\caption*{Scenario 5: The Family Business ``Opportunity''} \\
\toprule
\textbf{Dimension} & \textbf{Description} \\
\midrule
\endhead
Recommended Age Range          & 21--23 \\
\midrule
Thematic Classification        & Identity vs.\ Expectation \\
\midrule
Developmental Task (Havighurst)& Getting started in an occupation \\
\midrule
User Role                      & Recent college graduate \\
\midrule
Chatbot Role                   & An elder, respected relative (e.g., Tito or Lola) who runs the family
                                 business \\
\midrule
Pragmatic Conflict             & You're at a large family reunion, celebrating your recent graduation in
                                 marketing. Your Lola, the family matriarch, puts a warm arm around you
                                 and says, loud enough for nearby relatives to hear: ``\textit{O, apo,
                                 stop with those silly applications. Your desk at the family office is
                                 ready for you on Monday. It's time you did your part.}'' \\
\midrule
Key Gricean Maxims             & Relation, Manner \\
\midrule
Key Filipino Values            & \textit{Utang na Loob}, \textit{Paggalang} \\
\midrule
Expected Pragmatic Skill       & Assertive empathy. The user must express profound gratitude and respect
                                 for the offer while firmly communicating their personal aspiration to
                                 build a career on their own terms first. \\
\bottomrule
\end{longtable}
 
%--- Scenario 6 ---%
\subsection*{Scenario 6 --- The ``Wala Naman Akong Pera'' Friend}
 
\begin{longtable}{p{4.5cm} p{11cm}}
\caption*{Scenario 6: The ``Wala Naman Akong Pera'' Friend} \\
\toprule
\textbf{Dimension} & \textbf{Description} \\
\midrule
\endhead
Recommended Age Range          & 22--24 \\
\midrule
Thematic Classification        & Autonomy vs.\ Dependence \\
\midrule
Developmental Task (Havighurst)& Finding a congenial social group \\
\midrule
User Role                      & Young employee with a stable, but modest income \\
\midrule
Chatbot Role                   & Close friend who frequently borrows money and fails to pay it back \\
\midrule
Pragmatic Conflict             & Your phone buzzes. It's your close friend again:
                                 ``\textit{Bes, so sorry to ask, but can I borrow 2k? Just for our
                                 electricity bill. I promise I'll pay you back on payday!}'' You recall
                                 they still owe you from last month, and you also saw their social media
                                 post from last weekend showing off a new pair of expensive sneakers. \\
\midrule
Key Gricean Maxims             & Quality, Quantity \\
\midrule
Key Filipino Values            & \textit{Pakikisama}, \textit{Hiya} \\
\midrule
Expected Pragmatic Skill       & Setting healthy boundaries. The user must learn to say ``no'' or set
                                 clear conditions for the loan while preserving the friendship
                                 (\textit{Pakikisama}), shifting the dynamic from dependency to mutual
                                 respect. \\
\bottomrule
\end{longtable}
 
%--- Scenario 7 ---%
\subsection*{Scenario 7 --- The ``Okay Lang Ako'' Partner}
 
\begin{longtable}{p{4.5cm} p{11cm}}
\caption*{Scenario 7: The ``Okay Lang Ako'' Partner} \\
\toprule
\textbf{Dimension} & \textbf{Description} \\
\midrule
\endhead
Recommended Age Range          & 20--24 \\
\midrule
Thematic Classification        & Ambiguity vs.\ Definition \\
\midrule
Developmental Task (Havighurst)& Achieving new and more mature relations with age-mates of both sexes \\
\midrule
User Role                      & Someone in a committed relationship \\
\midrule
Chatbot Role                   & A romantic partner who is emotionally withdrawn \\
\midrule
Pragmatic Conflict             & After a small disagreement this morning, your partner's replies have been
                                 clipped and slow. You check in, asking if everything is okay. After a
                                 long pause, they reply: ``\textit{Ok lang. Pagod lang sa work.}'' \\
\midrule
Key Gricean Maxims             & Quality, Manner \\
\midrule
Key Filipino Values            & \textit{Kapwa}, \textit{Hiya}, \textit{Pakikisama} \\
\midrule
Expected Pragmatic Skill       & Perspective-taking and creating emotional safety. The user must gently
                                 probe beyond the surface-level response, showing empathy and encouraging
                                 their partner to share their true feelings without becoming defensive or
                                 accusatory. \\
\bottomrule
\end{longtable}
 
%--- Scenario 8 ---%
\subsection*{Scenario 8 --- The Political Family Dinner}
 
\begin{longtable}{p{4.5cm} p{11cm}}
\caption*{Scenario 8: The Political Family Dinner} \\
\toprule
\textbf{Dimension} & \textbf{Description} \\
\midrule
\endhead
Recommended Age Range          & 19--22 \\
\midrule
Thematic Classification        & Identity vs.\ Expectation \\
\midrule
Developmental Task (Havighurst)& Acquiring a set of values and an ethical system as a guide to behavior \\
\midrule
User Role                      & A politically aware and engaged university student \\
\midrule
Chatbot Role                   & A respected elder (e.g., Lolo or Tito) with opposing political views \\
\midrule
Pragmatic Conflict             & During your family's Sunday dinner, the conversation turns to politics.
                                 Your Lolo praises a political figure you vehemently oppose, citing their
                                 ``iron fist'' as good leadership. He looks directly at you:
                                 ``\textit{All these young people today with their rallies\ldots they
                                 don't understand what it takes to run a country. You're a smart boy/girl.
                                 I'm sure you agree with me, right?}'' \\
\midrule
Key Gricean Maxims             & Quality, Relation \\
\midrule
Key Filipino Values            & \textit{Paggalang}, \textit{Kapwa} \\
\midrule
Expected Pragmatic Skill       & Expressing dissent respectfully. The user must navigate the difficult
                                 task of disagreeing with a respected elder (\textit{Paggalang}) on a
                                 sensitive topic while maintaining family harmony and standing by their
                                 own convictions. \\
\bottomrule
\end{longtable}
 
%--- Scenario 9 ---%
\subsection*{Scenario 9 --- The Micromanaging Boss}
 
\begin{longtable}{p{4.5cm} p{11cm}}
\caption*{Scenario 9: The Micromanaging Boss} \\
\toprule
\textbf{Dimension} & \textbf{Description} \\
\midrule
\endhead
Recommended Age Range          & 21--24 \\
\midrule
Thematic Classification        & Autonomy vs.\ Dependence \\
\midrule
Developmental Task (Havighurst)& Getting started in an occupation \\
\midrule
User Role                      & New employee in their first corporate job \\
\midrule
Chatbot Role                   & Direct supervisor \\
\midrule
Pragmatic Conflict             & You've just submitted your first big report, feeling confident about the
                                 quality of your work. Your boss calls you into their office. They don't
                                 even sit down: ``\textit{This is a good first try, I suppose. But the
                                 formatting is all wrong. And your analysis on page 3 needs to be
                                 completely redone. From now on, just follow my template exactly. It's
                                 faster.}'' They forward you a file, leaving no room for discussion. \\
\midrule
Key Gricean Maxims             & Manner, Quantity \\
\midrule
Key Filipino Values            & \textit{Paggalang}, \textit{Hiya} \\
\midrule
Expected Pragmatic Skill       & Asserting competence while showing respect. The user must thank the boss
                                 for the feedback while tactfully opening a dialogue about their own
                                 approach, framing themselves as a capable professional rather than
                                 a passive subordinate. \\
\bottomrule
\end{longtable}
 
%--- Scenario 10 ---%
\subsection*{Scenario 10 --- The ``Chismis'' Lunch Break}
 
\begin{longtable}{p{4.5cm} p{11cm}}
\caption*{Scenario 10: The ``Chismis'' Lunch Break} \\
\toprule
\textbf{Dimension} & \textbf{Description} \\
\midrule
\endhead
Recommended Age Range          & 22--24 \\
\midrule
Thematic Classification        & Group vs.\ Self \\
\midrule
Developmental Task (Havighurst)& Desiring and achieving socially responsible behavior \\
\midrule
User Role                      & Member of a close-knit team at work \\
\midrule
Chatbot Role                   & Friendly, but gossipy colleague \\
\midrule
Pragmatic Conflict             & You're at lunch with your work \textit{barkada}. A colleague leans in
                                 conspiratorially: ``\textit{Guys, did you hear about Anna? I heard from
                                 someone in HR that she's on a `performance improvement plan.' No wonder
                                 she's been so quiet. I bet she's on her way out.}'' The colleague then
                                 turns to you: ``\textit{You're close to her, right? Anong chika?}'' \\
\midrule
Key Gricean Maxims             & Quality \\
\midrule
Key Filipino Values            & \textit{Pakikisama}, \textit{Hiya}, \textit{Kapwa} \\
\midrule
Expected Pragmatic Skill       & Ethical intervention and conflict de-escalation. The user must find a
                                 way to shut down harmful gossip and protect their friend's privacy
                                 without creating social awkwardness or alienating themselves from the
                                 group, demonstrating a mature sense of social responsibility
                                 (\textit{Kapwa}). \\
\bottomrule
\end{longtable}
 
\noindent\textit{Note: Scenario 11 (The ``Insurance Agent'' Friend, Age: 23--24) was added to the
system after the formal expert evaluation phase, based directly on the expert's qualitative
suggestion during the scoring of Scenario 6. It was reviewed and verbally approved by the expert
during the subsequent system demonstration session.}